% Problem record op_d754e0d170c01f86. This TeX file is derived from
% database/problems_json/op_d754e0d170c01f86.json by scripts/sync-tex.mjs; edit the JSON record
% and rerun the script rather than editing this file.
\section{Diamond-smoothed max-relative-entropy AEP for quantum channels}
\label{sec:op_d754e0d170c01f86}

\paragraph{Problem.}
Does the max-relative entropy of finite-dimensional quantum channels satisfy an
asymptotic equipartition property under uniform diamond-norm smoothing?  Let
$\mathcal N,\mathcal M:\mathcal L(A)\to\mathcal L(B)$ be quantum channels
with $D_{\max}(\mathcal N\|\mathcal M)<\infty$.  Using unnormalized Choi
operators, define the channel max-relative entropy and its smoothed version by
\begin{equation}
  \begin{aligned}
  D_{\max}(\mathcal N\|\mathcal M)
    &:=\inf\{\lambda:J_{\mathcal N}\leq2^\lambda J_{\mathcal M}\},\\
  D_{\max}^{\varepsilon}(\mathcal N\|\mathcal M)
    &:=\inf_{\substack{\widetilde{\mathcal N}\ {\rm CPTP}:\\
              \frac12\|\widetilde{\mathcal N}-\mathcal N\|_\diamond
              \leq\varepsilon}}
       D_{\max}(\widetilde{\mathcal N}\|\mathcal M).
  \end{aligned}
  \label{eq:p65-smooth-channel-max}
\end{equation}
The smoothing in Eq.~\eqref{eq:p65-smooth-channel-max} requires one channel
$\widetilde{\mathcal N}$ that approximates $\mathcal N$ uniformly over all
ancilla-assisted inputs.  Define
\begin{equation}
  \begin{aligned}
  D_{\max}^{\varepsilon,\infty}(\mathcal N\|\mathcal M)
    &:=\limsup_{n\to\infty}\frac1n
       D_{\max}^{\varepsilon}
       (\mathcal N^{\otimes n}\|\mathcal M^{\otimes n}),\\
  D_{\rm ch}^{\infty}(\mathcal N\|\mathcal M)
    &:=\lim_{n\to\infty}\frac1n
      \sup_{\psi_{R A^n}}
      D\!\left((\operatorname{id}_R\otimes\mathcal N^{\otimes n})(\psi)
      \middle\|
      (\operatorname{id}_R\otimes\mathcal M^{\otimes n})(\psi)\right),
  \end{aligned}
  \label{eq:p65-regularized-divergences}
\end{equation}
where $R\simeq A^{\otimes n}$ suffices and $D$ is quantum relative entropy.
Is the following identity valid for every such channel pair, and can the
$\limsup$ in Eq.~\eqref{eq:p65-regularized-divergences} be replaced by a
limit?
\begin{equation}
  \sup_{\varepsilon>0}
  D_{\max}^{\varepsilon,\infty}(\mathcal N\|\mathcal M)
  =D_{\rm ch}^{\infty}(\mathcal N\|\mathcal M).
  \label{eq:p65-channel-aep}
\end{equation}

\subsection*{Status}

Unsolved

\subsection*{Source}

Winter first proposed the identity in Eq.~\eqref{eq:p65-channel-aep}; Liu and
Winter discussed it formally in the setting of channel-resource erasure
\sourcecite{ref:p65-liu-winter}{LW19}.  Hirche restated it explicitly as the
technical conjecture needed in quantum-network discrimination
\sourcecite{ref:p65-hirche}{Hir23}.

\subsection*{Progress}

\begin{itemize}
  \item Gour and Winter proved asymptotic equipartition statements for two
  channel-resource divergences using a more permissive ``liberal'' smoothing.
  That smoothing does not require the single uniformly diamond-close channel
  in Eq.~\eqref{eq:p65-smooth-channel-max}, so it does not establish
  Eq.~\eqref{eq:p65-channel-aep}
  \sourcecite{ref:p65-gour-winter}{GW19}.

  \item Hirche proved the general bounds
  \begin{equation}
    D_{\rm ch}^{\infty}(\mathcal N\|\mathcal M)
    \leq D_{\max}^{\varepsilon,\infty}(\mathcal N\|\mathcal M)
    \leq D_{\max}(\mathcal N\|\mathcal M),
    \label{eq:p65-known-bounds}
  \end{equation}
  The bounds in Eq.~\eqref{eq:p65-known-bounds} do not identify the
  asymptotic rate.  Hirche also showed that Eq.~\eqref{eq:p65-channel-aep}
  would collapse the amortized
  Umegaki relative entropy of superchannels to their regularized relative
  entropy \sourcecite{ref:p65-hirche}{Hir23}.

  \item Fang, Gour, and Wang related an analogous AEP for unstabilized channel
  divergences without quantum-memory assistance to strong converses for
  channel discrimination.  Their latest formulation leaves that general
  strong-converse problem unresolved and does not prove the stabilized,
  diamond-smoothed identity in Eq.~\eqref{eq:p65-channel-aep}
  \sourcecite{ref:p65-fang-gour-wang}{FGW25}.

  \item Gour gives a counterexample in Theorems 11 and 15 of a 2026 preprint
  \sourcecite{ref:p65-gour-aep}{Gou26}.  For
  $0<t<(\sqrt5-2)^2$ and $0<\theta<2\arcsin(1/8)$, set
  \begin{equation}
    \begin{aligned}
      \mathcal N&=\operatorname{id}_3,\qquad
      \mathcal M=t\operatorname{id}_3+(1-t)\mathcal U,\\
      \mathcal U(X)&=UXU^\dagger,\qquad
      U=\operatorname{diag}(1,e^{i\theta},e^{-i\theta}).
    \end{aligned}
    \label{eq:p65-qutrit-counterexample}
  \end{equation}
  The channels in Eq.~\eqref{eq:p65-qutrit-counterexample} have finite
  $D_{\max}$.  For every fixed $0<\varepsilon<\sin^2\theta$,
  \begin{equation}
    \begin{aligned}
      \lim_{n\to\infty}\frac1nD_{\max}^{\varepsilon}
        (\mathcal N^{\otimes n}\|\mathcal M^{\otimes n})
        &=\log_2(1/t),\\
      D_{\rm ch}^{\infty}(\mathcal N\|\mathcal M)
        &\leq\tfrac12\log_2(1/t)+\log_2\frac4{1-t}
        <\log_2(1/t).
    \end{aligned}
    \label{eq:p65-aep-separation}
  \end{equation}
  Monotonicity and the unsmoothed upper bound make the supremum in
  Eq.~\eqref{eq:p65-channel-aep} equal to $\log_2(1/t)$.
  Equation~\eqref{eq:p65-aep-separation} therefore refutes the proposed
  identity, including its vanishing-error version (Section VII, Eqs.~(232)--(234)).
\end{itemize}

\subsection*{References}

\begin{enumerate}[label={},leftmargin=4.8em,labelsep=0.5em]
  \footnotesize
  \item[\textup{[LW19]}]\label{ref:p65-liu-winter}
  Z.-W. Liu and A. Winter,
  ``Resource Theories of Quantum Channels and the Universal Role of Resource
  Erasure,'' arXiv preprint (2019).
  \href{https://arxiv.org/abs/1904.04201}{arXiv:1904.04201}.

  \item[\textup{[GW19]}]\label{ref:p65-gour-winter}
  G. Gour and A. Winter, ``How to Quantify a Dynamical Quantum Resource,''
  \emph{Physical Review Letters} \textbf{123}, 150401 (2019).
  \href{https://doi.org/10.1103/PhysRevLett.123.150401}{doi:10.1103/PhysRevLett.123.150401};
  \href{https://arxiv.org/abs/1906.03517}{arXiv:1906.03517}.

  \item[\textup{[Hir23]}]\label{ref:p65-hirche}
  C. Hirche, ``Quantum Network Discrimination,''
  \emph{Quantum} \textbf{7}, 1064 (2023).
  \href{https://doi.org/10.22331/q-2023-07-25-1064}{doi:10.22331/q-2023-07-25-1064};
  \href{https://arxiv.org/abs/2103.02404}{arXiv:2103.02404}.

  \item[\textup{[FGW25]}]\label{ref:p65-fang-gour-wang}
  K. Fang, G. Gour, and X. Wang,
  ``Towards the Ultimate Limits of Quantum Channel Discrimination and Quantum
  Communication,'' \emph{Science China Information Sciences} \textbf{68},
  180509 (2025).
  \href{https://doi.org/10.1007/s11432-024-4488-0}{doi:10.1007/s11432-024-4488-0};
  \href{https://arxiv.org/abs/2110.14842}{arXiv:2110.14842}.

  \item[\textup{[Gou26]}]\label{ref:p65-gour-aep}
  G. Gour, ``Failure of the Asymptotic Equipartition Property for Quantum
  Channels,'' arXiv preprint (2026).
  \href{https://arxiv.org/abs/2609.15004v1}{arXiv:2609.15004v1}.
\end{enumerate}

\subsection*{Comment}

The universal identity in Eq.~\eqref{eq:p65-channel-aep} is false by
Gour's preprint \sourcecite{ref:p65-gour-aep}{Gou26}.  The limit exists for
this counterexample family; the separate question of its existence for
arbitrary channel pairs is not settled by this result.  The record therefore
remains Unsolved while both requests are retained in the statement.  The
counterexample also does not decide
\href{https://qiqc-op.com/problem/op_08387c140f552732/}{the fixed-error parallel channel Stein problem}
or the conditional equality in
\href{https://qiqc-op.com/problem/op_a4600b38b94042a8/}{the parallel versus nested-adaptive superchannel-discrimination problem}.

\subsection*{Field}

Quantum Communication

\subsection*{Topic}

Quantum relative entropy; One-shot and finite-blocklength bounds

\subsection*{ID}

\texttt{op\_d754e0d170c01f86}
